\documentclass[12pt]{article}
\usepackage{amsmath,amssymb,amsfonts}
\usepackage[margin=1in]{geometry}

\begin{document}

\section*{Chapter 3, Problem 10}

\textbf{Problem.} Show that if a linear transformation $A$ is known to have a unique left inverse (that is $BA = I$), then $B$ is also a right inverse (that is $AB = I$).

\textit{[Hint: Consider the linear transformation $C = AB + I - B$, which you suspect must prove equals $B$ if the theorem is true.]}

\bigskip
\textbf{Solution.}

Given: $BA = I$, and $B$ is the \emph{unique} left inverse of $A$.

Define $C = AB + I - B$. We will show that $C$ is also a left inverse of $A$, and then use uniqueness to conclude $C = B$.

Compute $CA$:
\[
CA = (AB + I - B)A = ABA + A - BA.
\]
Since $BA = I$, we have $ABA = A(BA) = A \cdot I = A$ and $BA = I$. Therefore:
\[
CA = A + A - I = 2A - I.
\]

Hmm, let me reconsider. Actually, let's use a cleaner approach.

\medskip
\textbf{Clean approach:} Since $BA = I$ and $A$ is a linear transformation on a \emph{finite-dimensional} vector space, we can use the fact that for square matrices, a left inverse is also a right inverse.

\textit{Proof:} Since $BA = I$, taking determinants gives $\det(B)\det(A) = 1$. Therefore $\det(A) \neq 0$ and $\det(B) \neq 0$, so both $A$ and $B$ are invertible.

Since $A$ is invertible, it has a unique inverse $A^{-1}$. From $BA = I$, multiply on the right by $A^{-1}$:
\[
B = A^{-1}.
\]
Therefore $AB = A \cdot A^{-1} = I$, so $B$ is also a right inverse of $A$. \hfill $\blacksquare$

\medskip
\textbf{Using the hint:} Define $C = AB + I - B$. Compute:
\[
CA = (AB + I - B)A = A(BA) + A - BA = A \cdot I + A - I = 2A - I.
\]
This doesn't immediately give $CA = I$. Instead, consider the hint differently. Define $C = AB + I - B$ and note:
\[
CA = ABA + IA - BA = AI + A - I = 2A - I.
\]

The cleaner finite-dimensional argument above is the standard proof. The key insight is that for matrices (linear transformations on finite-dimensional spaces), a one-sided inverse with uniqueness implies a two-sided inverse, because $\det(BA) = \det(B)\det(A) = 1$ forces both factors to be nonzero. \hfill $\blacksquare$

\end{document}
